\begin{helper}[Restricted maximality gives terminal minus-closedness]
Let $K_J$ be a good hierarchy on $J$ and $K_U$ a good hierarchy on $U$.
Suppose their last Gamma pairs agree, the terminal stored set of $K_U$ is
contained in the last indexed stored set of $K_J$, and every last Gamma
pair of $K_J$ avoids the restricted segment colors at $U$.  If no
one-step Tashkinov augmentation of $U$ preserves all those avoidance
conditions, then $U$ is minus-closed outside the union of the last Gamma
pairs indexed by the terminal stored set of $K_U$.  Here one also assumes
the companion comparison supplied by the common computed stored set: a
color in the last indexed set of $K_J$ but outside the terminal set of
$K_U$ is missing on $U$.
\end{helper}
\begin{proof}
Let $D$ be the last indexed stored set of $K_J$, let $D_U$ be the
terminal stored set of $K_U$, and write $\Gamma_\eta$ for the common last
Gamma pair.  Suppose the asserted
\noderef{TreeColorSetMinusClosed} property fails.  Unfolding it and
\noderef{TreeColorSetClosed} gives a color $c$ which is missing on $U$,
does not lie in $\bigcup_{\eta\in D_U}\Gamma_\eta$, and colors an exact
boundary-edge identity $f$ of $U$.

Apply \noderef{BoundaryMissingColorAugmentation} to $c$ and $f$.  It
produces a tree $W$ obtained by adjoining exactly $f$, together with a
one-step $\sigma$-\noderef{TashkinovAugmentation} from $U$ to $W$.  Thus
$U\preccurlyeq W$, $|W|=|U|+1$, every old entry is unchanged, the new
last edge is $f$, and its color is $c$.

Fix $\eta\in D$.  First suppose $\eta\in D_U$.  Since $c$ is outside the
union indexed by $D_U$ and the last Gamma pairs agree, we have
$c\notin\Gamma_\eta$.  Apply
\noderef{HierarchyFirstMissingPositionPrefixCutoff} to $U\preccurlyeq W$
with cutoff $|U|$.  For an edge position contributing to the restricted
segment of $W$, every position below $|U|$ has the same edge, color, and
strict first-missing cutoff as in $U$.  The only other possible position
is the new position $|U|$, whose color is $c$.  Hence the restricted
color set at $W$ is contained in the corresponding set at $U$ together
with $\{c\}$.  Endpoint avoidance and $c\notin\Gamma_\eta$ give the
required avoidance at $W$.  The cutoff theorem includes its no-missing
fallback, so this argument does not assume an unprovided missing witness.

Now suppose $\eta\notin D_U$.  Because $\eta\in D$ and $D_U$ is obtained
outside the terminal set, the assumed companion comparison gives
$\eta\in\noderef{TreeMissingColors}(\sigma,U)$, so $\eta$ is missing at
an old vertex of $U$.
All old vertex entries are unchanged in $W$.  Therefore the least old
vertex missing $\eta$, and the absence of such a witness at every earlier
position, show directly from \noderef{HierarchyFirstMissingPosition} that
the first missing position in $U$ and $W$ is the same.  The new edge
position $|U|$ is not strictly before this old position, and all earlier
edge entries agree.  Thus the restricted segment-color sets at $U$ and
$W$ are equal, so endpoint avoidance again persists.

We have proved the fixed-Gamma avoidance condition at $W$ for every
$\eta\in D$.  Together with the one-step augmentation, this contradicts
the assumed endpoint maximality.  Hence no such $c$ and $f$ exist, which
is exactly the desired terminal minus-closedness.
\end{proof}
